% !TEX root = ../main.tex

\chapter{Background and Related Work}
\label{chap:background}

\section{LLMs in Mental Health Applications}

\todo{Review digital and AI-assisted mental-health support, from established chatbot systems to recent open-ended LLM use, with emphasis on safety limitations.}

\subsection{Existing Digital Mental Health Chatbots}

\todo{Draft subsection.}

\subsection{LLM-Based Mental Health Support}

\todo{Draft subsection.}

\subsection{Safety Limitations of Existing Approaches}

\todo{Draft subsection.}

\section{Neuro-Symbolic Approaches to AI Safety}

\todo{Introduce neuro-symbolic AI as a methodological approach for combining flexible neural language understanding with explicit symbolic structures.}

\subsection{Definition and Basic Idea}

\todo{Draft subsection.}

\subsection{Benefits for Safety-Critical Systems}

\todo{Draft subsection.}

\subsection{Relevance for Clinical LLM Safety}

\todo{Draft subsection.}

\section{Imagery Rehearsal Therapy as the Clinical Use Case}

\todo{Explain IRT as the therapeutic protocol targeted by the thesis, focusing on therapeutic rationale, protocol structure, and the rescripting stage.}

\subsection{Treatment Rationale of IRT}

\todo{Draft subsection.}

\subsection{Structure of the IRT Protocol}

\todo{Draft subsection.}

\subsection{The Rescripting Stage as the Safety-Critical Point}

\todo{Draft subsection.}

\begin{figure}
  \centering
  \fbox{\parbox{0.85\linewidth}{\centering Figure placeholder: IRT protocol flow with rescripting highlighted.}}
  \caption{IRT protocol flow with the rescripting stage highlighted as the safety-critical interaction point.}
  \label{fig:irt-protocol-flow}
\end{figure}

\section{Adaptive and Maladaptive Rescripting Patterns as Clinical Safety Categories}

\todo{Define the clinical categories used by the safety layer. Adaptive categories are grounded in mastery-oriented IRT and MMS concepts; maladaptive categories are derived from IRT-relevant clinical failure modes.}

\subsection{Adaptive Rescripting Patterns}

\todo{Draft subsection.}

\subsection{Maladaptive Rescripting Patterns}

\todo{Draft subsection.}

\subsection{From Clinical Patterns to Computational Categories}

\todo{Draft subsection.}

\begin{table}
  \centering
  \caption{Adaptive and maladaptive rescripting patterns with definitions, examples, and clinical rationale/source basis.}
  \label{tab:clinical-safety-categories}
  \fbox{\parbox{0.85\linewidth}{\centering Table placeholder.}}
\end{table}

\begin{figure}
  \centering
  \fbox{\parbox{0.85\linewidth}{\centering Figure placeholder: neuro-symbolic framing of the safety layer.}}
  \caption{Neuro-symbolic framing of the proposed safety layer. The LLM component maps free-form rescripting text to clinical safety categories, while the symbolic ontology provides explicit category definitions, severity levels, and correction relations.}
  \label{fig:neuro-symbolic-framing}
\end{figure}
